\section*{Capítulo \ref{ch:b2:diffraction} --- Difração de Fraunhofer e filtragem espacial}

\begin{solution}{exo:b2:diffraction:1}
$2\lambda L/b = \qty{25}{mm}$; \omterm{thm:b2:gratings:nwaves}{máximos secundários} em $1.43\lambda L/b = \qty{18}{mm}$ e
$2.46\lambda L/b = \qty{31}{mm}$; \omterm{rem:b2:diffraction:fresnelnumber}{número de Fresnel} $b^2/\lambda L = 8 \times 10^{-3}$: longe
o bastante. Com a lente, $2\lambda f/b = \qty{6.3}{mm}$.
\end{solution}

\begin{solution}{exo:b2:diffraction:2}
$1.22\lambda/D$: olho \qty{2.2e-4}{rad} (\qty{46}{''}); \qty{10}{cm}: \qty{1.4}{''}; Hubble
\qty{0.058}{''}; antena $1.22 \times 0.06/64 = \qty{1.1e-3}{rad} = \qty{3.9}{'}$; objetiva:
$\qty{1.1e-4}{rad} \times \qty{50}{mm} = \qty{5.6}{\micro m}$ de raio, um disco de \qty{11}{\micro m} ---
três pixels: a \omterm{def:b2:diffraction:huygens}{difração} já limita um sensor a $f/8$.
\end{solution}

\begin{solution}{exo:b2:diffraction:3}
$\lambda f/a = \qty{2}{mm}$; envoltória $2\lambda f/b = \qty{20}{mm}$; dez franjas; as ordens
$p = 5, 10, \dots$ ($a/b = 5$) estão ausentes.
\end{solution}

\begin{solution}{exo:b2:diffraction:4}
$\theta = 1.22\lambda/D = \qty{3.9e-4}{rad}$: uma mancha de \qty{290}{km} de diâmetro; com \qty{3.5}{m},
\qty{170}{m}. A própria \omterm{def:b2:diffraction:huygens}{difração} do refletor, $\lambda/\qty{1}{m}$, espalha o
retorno por \qty{600}{m}: dos $10^{17}$ fótons enviados, voltam uns poucos.
\end{solution}

\begin{solution}{exo:b2:diffraction:5}
(a) $\int_{-b/2}^{b/2}\eu^{\iu kx\sin\theta}\dd x = b\,\operatorname{sinc}(kb\sin\theta/2)$. (b)
$\dd(\sin x/x)/\dd x = 0$ dá $\tan x = x$; $x = 4.49 = 1.43\pi$, $\operatorname{sinc}^2 =
0.047$. (c) $0.903$. (d) A fase se torna $kx(\sin\theta - \sin\theta_0)$: o
mesmo padrão, em torno de $\theta_0$.
\end{solution}

\begin{solution}{exo:b2:diffraction:6}
$0.61\lambda/\mathrm{NA}$: \qty{220}{nm}; \qty{170}{nm}; \qty{110}{nm}. Além de uns $500\times$
a ocular só amplia o borrão (ampliação vazia). Elétrons
a \qty{4}{pm} com $\mathrm{NA} \approx 0.01$: \qty{0.2}{nm} --- os átomos.
\end{solution}

\begin{solution}{exo:b2:diffraction:7}
(a) A integral se fatora em $x$ e $y$. (b) $2\lambda f/b = \qty{5}{mm}$. (c) As
cordas de um disco são mais curtas que seu diâmetro fora do centro: sua
largura efetiva é menor e seu padrão, mais largo. (d) Quatro espigões de
\omterm{def:b2:diffraction:huygens}{difração} em cruz.
\end{solution}

\begin{solution}{exo:b2:diffraction:8}
(a) $\sin\theta = \pm0.275$ (\qty{16}{\degree}), $\pm0.55$ (\qty{33}{\degree}). (b) $\sin\alpha \ge 0.275$.
(c) Apenas as ordens $0$, $\pm1$: uma imagem senoidal de período correto, e não
as barras nítidas. (d) $\sin\alpha \ge \lambda/2d = 0.14$: o dobro da resolução.
\end{solution}

\begin{solution}{exo:b2:diffraction:9}
(a) $2.44\lambda f/D = \qty{15}{\micro m}$; furo de \qty{23}{\micro m}. (b) $\lambda/\qty{50}{\micro m} =
\qty{1.3e-2}{rad}$: \qty{0.25}{mm} fora do eixo, bloqueada. (c) $1\%$; o furo guarda só
as \omterm{prop:b2:diffraction:fourier}{frequências espaciais} baixas. (d) Descentrado de alguns micrômetros, o
furo corta a própria mancha.
\end{solution}

\begin{solution}{exo:b2:diffraction:10}
(a) As amplitudes somam a onda sem obstáculo, que é nula fora do eixo:
$|\underline s_{\text{fenda}}| = |\underline s_{\text{tira}}|$ ali. (b) Franjas separadas por $\lambda L/b =
\qty{16}{mm}$; meça o espaçamento e obtenha $b$. (c) A sombra da tira fica
dentro do feixe brilhante não difratado. (d) Gotículas de diâmetro $d$ difratam
como furos: anéis em $1.22\lambda/d$, com o vermelho por fora; o raio dá $d$.
\end{solution}

\begin{solution}{exo:b2:diffraction:11}
(a) $|1 + \iu\varphi|^2 = 1 + \varphi^2$: uniforme em primeira ordem. (b) Sem a
luz axial, $|\iu\varphi|^2 = \varphi^2$: segunda ordem, fraca. (c) $|\iu + \iu\varphi|^2 =
1 + 2\varphi + \dots$: linear. (d) A mancha axial ofusca a luz
difratada; atenuá-la iguala as duas amplitudes e aumenta o contraste.
\end{solution}

\begin{solution}{exo:b2:diffraction:12}
(a) Semilargura $\lambda/W$; dividida por $p/a\cos\theta$: $\Delta\lambda = \lambda a\cos\theta/pW \approx
\lambda/pN$. (b) O feixe colimado tem diâmetro $Df_{\text{col}}/f_{\text{tel}}$;
uma rede menor desperdiçaria luz e poder de resolução. (c) Uma
resolução é sempre o ângulo ao longo do qual a fase através da
abertura varia de um comprimento de onda, $\lambda/(\text{tamanho})$. (d) A largura de
sua imagem, que domina quando é maior que a largura de \omterm{def:b2:diffraction:huygens}{difração} da
rede --- o caso habitual.
\end{solution}

\begin{solution}{pb:b2:diffraction:1}
\textbf{1.} $1.22\lambda/D = \qty{2.2e-4}{rad}$; $\times\qty{17}{mm} = \qty{3.8}{\micro m}$: dois
espaçamentos de cones --- a retina amostra no limite da óptica.

\textbf{2.} \qty{56}{\micro m} a \qty{25}{cm}, \qty{2.2}{mm} a \qty{10}{m}; \qty{1}{mm} a
\qty{50}{cm} é dez vezes o limite: com folga.

\textbf{3.} \qty{9.6e-5}{rad}; as aberrações do cristalino em plena abertura e os
bastonetes, mais grosseiros e agrupados, anulam o ganho.

\textbf{4.} $1.5/2.2 \times 10^{-4} = \qty{7}{km}$.

\textbf{5.} A \omterm{def:b2:diffraction:huygens}{difração} piora para \qty{6.7e-4}{rad}, mas o círculo de borrão de
um olho desfocado encolhe com a abertura: a profundidade de campo vence.

\textbf{6.} \qty{6.7e-7}{rad} (\qty{0.14}{''}); $\times\qty{10}{m} = \qty{6.7}{\micro m}$ de raio: um
disco de $1.3$ pixel.

\textbf{7.} \qty{1}{''} $= \qty{4.9e-6}{rad}$, sete vezes o limite; $D = 1.22\lambda/4.9
\times 10^{-6} = \qty{14}{cm}$ já alcança o seeing.

\textbf{8.} \qty{2.7e-6}{rad} (\qty{0.55}{''}); os erros de fase da atmosfera são
uma fração menor de um comprimento de onda maior: correções menos numerosas e mais lentas.

\textbf{9.} $0.15 < 0.55$: não; em \qty{550}{nm} a partir do espaço, \qty{0.14}{''}: no limite.

\textbf{10.} $6.7 \times 10^{-7} \times 3.8 \times 10^8 = \qty{250}{m}$; o módulo lunar, não.

\textbf{11.} Uma cruz de quatro espigões de \omterm{def:b2:diffraction:huygens}{difração}.

\textbf{12.} $(1000/7)^2 = 2 \times 10^4$: objetos mais fracos, não mais nítidos.

\textbf{13.} \qty{370}{nm} e \qty{240}{nm}: $2700$ e $4200$ pares de linhas por
milímetro.

\textbf{14.} A bactéria, sim; sua estrutura de \qty{0.2}{\micro m} no limite
com óleo, não a seco; o vírus, não.

\textbf{15.} $d \ge \lambda/\mathrm{NA} = \qty{390}{nm}$ sob iluminação normal, \qty{200}{nm}
com luz oblíqua --- o $0.61\lambda/\mathrm{NA}$ de dois pontos fica entre os dois.

\textbf{16.} O índice do óleo aumenta $n\sin\alpha$ e suprime a reflexão
total que prenderia os raios muito inclinados na lamínula.

\textbf{17.} \qty{110}{nm}; elétrons: $0.61 \times 4\,\text{pm}/0.01 = \qty{0.24}{nm}$.

\textbf{18.} A lâmina atrasa de um quarto de onda a luz não difratada
para que ela interfira com a luz difratada pela estrutura de fase da
célula; o padrão de fase se torna um padrão de intensidade.

\textbf{19.} \qty{4.3e-4}{rad}; $1.5 + 2 \times 43 = \qty{88}{mm}$.

\textbf{20.} $2.44\lambda f_1/D = \qty{8.7}{\micro m}$; furo de \qty{13}{\micro m}.

\textbf{21.} $f_2 = 20f_1 = \qty{200}{mm}$; \qty{2.2e-5}{rad}; $30 + 44 = \qty{74}{mm}$ depois de um
quilômetro.

\textbf{22.} $\lambda/\qty{100}{\micro m} = \qty{5.3e-3}{rad}$: \qty{53}{\micro m} fora do eixo no
plano do furo --- bloqueada.

\textbf{23.} $1\%$, absorvido pela placa do furo.

\textbf{24.} Um feixe de diâmetro $D$ dobra ao longo de $\sim D^2/\lambda$: \qty{4}{m} para
\qty{1.5}{mm}, \qty{1.7}{km} para \qty{30}{mm}.

\textbf{25.} Olho \qty{3}{mm}, telescópio \qty{1}{m}, objetiva (NA $1.4$), feixe
\qty{30}{mm}: em cada caso, $\lambda/D$ fixa o menor ângulo, o detalhe mais fino, o
espalhamento mais lento.
\end{solution}
